\begin{abstract}
The ARC WhiteBox Estimation Challenge asks for estimates of post-ReLU mean activations in random MLPs under Gaussian input, with a score that rewards both low final-layer error and low effective compute. This write-up describes a structure-aware estimator built around three observations. First, a diagonal Gaussian propagation already identifies many neurons as effectively dead, effectively linear, or near the ReLU kink. Second, expensive samples are most valuable on the remaining kink paths and on late-layer classifications whose sign errors propagate to the scored layer. Third, the ReLU masks induced by those samples create row-sparse matrix products that can be evaluated exactly without multiplying by known zeros. The resulting method combines analytic moment propagation, antithetic \sobol{} sampling, staged active-set refinement, and exact sparse/dense routing of sampled activations. On the public Phase 1 evaluator, the submitted artifact achieved adjusted final-layer score $9.121856491923\times 10^{-8}$, raw final-layer MSE $3.724918391868\times 10^{-7}$, mean multiplier $0.2416126721$, mean effective compute $65.719$G, and no public failures.
\end{abstract}